% !PRINTED 444
\weberpage{444}
%<<CONT
\gesperrt{1. Es giebt 24 siebenz\"ahlige Pole, die sich, den acht siebengliedrigen Cykeln entsprechend, in acht Tripel theilen. Jeder dieser 24 Punkte kann in jeden anderen durch Substitutionen der Gruppe $G_{168}$ transformirt werden.}

\begin{center}
\S.\,118.

\gesperrt{Die Hauptaxen.}
\end{center}

Wir gehen zur Betrachtung der Substitutionen $4^{\text{ten}}$ und $2^{\text{ten}}$ Grades \"uber, als deren Repr\"asentanten wir $\Theta$ und $\Theta^2$ w\"ahlen. Nach \S.\,115, (10) und \S.\,116 haben wir f\"ur $\Theta$ die cubische Gleichung zu l\"osen:
\weq{1}{
\begin{vmatrix}
\gamma\varepsilon^2-\lambda, & \alpha\varepsilon^6, & \beta\\
\alpha, & \beta\varepsilon^4-\lambda, & \gamma\varepsilon^5\\
\beta\varepsilon^3, & \gamma, & \alpha\varepsilon-\lambda
\end{vmatrix} = 0,
}
die entwickelt so lautet:
\weq{2}{
\begin{aligned}
&\lambda^3-\lambda^2(\alpha\varepsilon+\beta\varepsilon^4+\gamma\varepsilon^2)\\
&\quad -\lambda\left[\varepsilon^6(\alpha^2-\beta\gamma)+\varepsilon^3(\beta^2-\gamma\alpha)+\varepsilon^5(\gamma^2-\alpha\beta)\right]\\
&\quad +\alpha^3+\beta^3+\gamma^3-3\alpha\beta\gamma=0.
\end{aligned}
}
Diese Gleichung reducirt sich aber mit Benutzung der Relationen (8), (9), (14) des \S.\,115 auf
\weq{3}{\lambda^3-\lambda^2+\lambda-1=0,}
und hat die Wurzeln
\[
\lambda=1,\quad \lambda=\pm i,
\]
und daraus lassen sich die Coordinaten der drei Pole berechnen, die sich als rationale Functionen von $\varepsilon$ und $i$ ergeben.

F\"ur die Pole von $\Theta^2$ erhalten wir nach \S.\,115, (16) zun\"achst die cubische Gleichung
\[
\begin{vmatrix}
\beta-\lambda, & \gamma\varepsilon^3, & \alpha\varepsilon^2\\
\gamma\varepsilon^{-3}, & \alpha-\lambda, & \beta\varepsilon^{-1}\\
\alpha\varepsilon^{-2}, & \beta\varepsilon, & \gamma-\lambda
\end{vmatrix} = 0,
\]
die nach den Relationen zwischen $\alpha,\beta,\gamma$ leicht in die Form
\weq{4}{(\lambda+1)^2(\lambda-1)=0}
gebracht wird und daher eine Doppelwurzel $\lambda=-1$ und eine einfache Wurzel $\lambda=1$ hat. Setzen wir den Werth $\lambda=-1$
%CONT>>
